%×××××××××××××××××××××××××××××××
%中文封面（专硕）
%×××××××××××××××××××××××××××××××

\thispagestyle{empty}
\noindent \wuhao% 五号字号
{密级：\hfill 中图分类号：TP3x}

\noindent $\square$全日制   \qquad $\square$非全日制
\vspace{20pt}
\begin{figure}[h]%加入校徽
    \centering
    \includegraphics[width=0.9\textwidth]{./figures/zjgsulogo.png}
\end{figure}
    
\vspace{20pt}
\begin{center}
    {\chuhao \textbf{硕士学位论文}}
\end{center}
\vspace{20pt}
\begin{center}
    {\yihao \textbf{(专业学位)}}
\end{center}

\vspace{40pt}
\begin{center}
    \begin{tabular}{lp{10cm}<{\centering}} 
        {\sanhao \heiti \textbf{论文题目：}} & {\sanhao \heiti \textbf{居中}} \\\cline{2-2} \\% 使用&分隔两列，\\开始新行
        %& {\sanhao \heiti  如果有副标题去掉最前面的百分号%，然后替代这段话 }  \\\cline{2-2}\\
    \end{tabular} 
\end{center}
    
    
\vspace{60pt}
\begin{center}
    \begin{tabular}{lp{6cm}<{\centering}}   
        {\sanhao \kaishu \textbf{作者姓名：}} & {\sanhao \kaishu \textbf{浙小商}} \\\cline{2-2}\\ % 使用&分隔两列，\\开始新行，cline命令用来画下划线  

        {\sanhao \kaishu \textbf{专业学位类型：}} & {\sanhao \kaishu \textbf{电子信息}} \\\cline{2-2}\\ 

        {\sanhao \kaishu \textbf{专业学位领域：}} & {\sanhao \kaishu \textbf{计算机技术}} \\\cline{2-2}\\ 

        {\sanhao \kaishu \textbf{研究方向：}} & {\sanhao \kaishu \textbf{应用密码学}} \\\cline{2-2}\\

        {\sanhao \kaishu \textbf{指导教师：}} & {\sanhao \kaishu \textbf{李四、王五}}\\\cline{2-2}    
    \end{tabular} 
\end{center}

\vspace{100pt}
\begin{center}
    {\sanhao \kaishu \textbf{提交日期：20xx年xx月}}
\end{center}
